% JAC1964-20261121. /root/programme_audit; GPT-6 Astra Pro.
% Effort setting not exposed. New deductions with their exact scope.

\subsection*{JAC-P11. A positive CM orbit beyond the Picard summand}

\begin{proposition}[JAC-P11]\label{jac:P11}
Let $J$ be the Fermat-quartic Jacobian with symmetric principal theta
class $\theta$, and let $\eta=\mathcal F([C]_{(1)})$ be the
hyperflex-normalized class of AR-8. After a finite extension of the
number field $K$, there is a rational correspondence $T$ on
$\operatorname{CH}^2_{(1)}(J)_{\mathbf Q}$ such that
\[
 T^2=-1,\qquad h_K(Tz,Tw)=h_K(z,w),\qquad
 T\eta\in\operatorname{CH}^2(N).                       \tag{JAC.P11.1}
\]
Here $N$ is the same orthogonal complementary summand chosen in AR-8.
The classes $\eta,T\eta$ are independent even modulo algebraic
equivalence. For $r,s\in\mathbf Q$, not both zero, one has
\[
 h_K(r\eta+sT\eta,r\eta+sT\eta)
   =(r^2+s^2)h_K(\eta,\eta)>\frac{r^2+s^2}{100},        \tag{JAC.P11.2}
\]
where the final bound uses AR-13. Hence the height is strictly
positive on
\[
 \operatorname{CH}^2(M)(K)_{\mathbf Q}
       \oplus\mathbf Q\eta\oplus\mathbf Q T\eta.        \tag{JAC.P11.3}
\]
The last two summands form a one-dimensional $\mathbf Q(i)$-module;
no assertion that they exhaust $\operatorname{CH}^2(N)$ is made.
\end{proposition}

\begin{proof}
Bloch's isogeny gives $J\sim E^3$ geometrically, with $j(E)=1728$.
Thus the centre of $\operatorname{End}^0(J)$ is $\mathbf Q(i)$.
Choose its element $j$ with $j^2=-1$, and enlarge $K$ so it is
defined. Rosati preserves the centre and sends $j$ to $\pm j$.
If $j^\dagger=j$, then
$\operatorname{Tr}(j j^\dagger)=\operatorname{Tr}(-1)=-6$,
contrary to Rosati positivity. Hence $j^\dagger=-j$.
The precise background is Milne, \emph{Abelian Varieties}, v2.0,
Proposition~14.2 and Theorem~14.3, printed pp.~61--62.

Choose a positive integer $n$ such that $a=nj$ is an endomorphism.
It is an isogeny, and
\[
 a^2=[-n^2],\qquad a^\dagger a=[n^2],\qquad
 \deg(a)=n^6,\qquad a^*\theta=n^2\theta
     \text{ in }\operatorname{CH}^1(J)_{\mathbf Q}.     \tag{JAC.P11.4}
\]
For the last assertion the polarization identity first gives
equality in $\operatorname{NS}(J)_{\mathbf Q}$. The difference is
symmetric and algebraically trivial, hence two-torsion, so the
equality holds in the rational Chow group.
Set
\[
 T=n^{-3}a^*\bigm|_{\operatorname{CH}^2_{(1)}(J)_{\mathbf Q}}.
                                                        \tag{JAC.P11.5}
\]
Replacing $(n,a)$ by $(kn,[k]a)$ multiplies this pullback by
$k^{2\cdot2-1}=k^3$, so the definition is independent of clearing
denominators. Since $[-n^2]^*$ acts on this grade by
$(-n^2)^3$, we get $T^2=-1$ exactly in rational Chow groups.
The height projection formula gives
\[
 h_K(a^*z,a^*w)=h_K(z,a_*a^*w)=\deg(a)h_K(z,w)
              =n^6h_K(z,w).
\]
Thus $T$ is an isometry. It preserves primitiveness because
$\theta a^*z=n^{-2}a^*(\theta z)$.

Let $G=\mu_4^2$ and $p_G$ be the averaging projector of AR-8.
Centrality of $a$ makes $T$ commute with $G$, so $T\eta$ is
$G$-invariant. AR-8 constructs its Picard projector $e'$ with
the exact relation $e'p_G=0$. Therefore $e'T\eta=0$; since $T\eta$
is primitive, it belongs to that same $N$.
Symmetry and isometry give
\[
 h_K(\eta,T\eta)=h_K(T\eta,T^2\eta)
               =-h_K(T\eta,\eta)=-h_K(\eta,T\eta).
\]
The mixed term is zero, and both diagonal terms equal
$h_K(\eta,\eta)>1/100$ by AR-13.

Independence modulo algebraic equivalence can also be proved
without using the height. If $r\eta+sT\eta$ is algebraically
trivial, apply $r-sT$. Algebraic correspondences preserve
algebraic equivalence, so $(r^2+s^2)\eta$ is algebraically trivial.
Bloch's infinite-order theorem, transported by Fourier as in AR-8,
forces $r=s=0$. Finally AR-2 gives positivity on
$\operatorname{CH}^2(M)$ and height orthogonality to $N$,
which proves (JAC.P11.3).
\end{proof}

\subsection*{JAC-P12. The exact scope of tautological generation}

\begin{proposition}[JAC-P12]\label{jac:P12}
Let $R(J)$ be the usual tautological ring of the hyperflex-embedded
Fermat quartic: the smallest intersection subring of
$\operatorname{CH}^*(J)_{\mathbf Q}$ containing $[C]$ and stable
under Fourier and all $[n]^*$. Then
\[
 R^2(J)=\mathbf Q\theta^2\oplus\mathbf Q\eta,\qquad
 R^2(J)_{\mathrm{hom}}=\mathbf Q\eta.                  \tag{JAC.P12.1}
\]
The class $T\eta$ from JAC-P11 does not belong to $R^2(J)$ even
after adding an algebraically trivial class. Thus the ordinary
tautological generation theorem cannot identify the full primitive
grade-one Chow group with the Fermat Ceresa line.
\end{proposition}

\begin{proof}
Polishchuk, \emph{Lie symmetries of the Chow group of a Jacobian
and the tautological subring}, arXiv:math/0504448v3,
Theorem~0.2(i), p.~2, proves that $R(J)$ is generated by
$p_n=\mathcal F([C]_{(n-1)})$ and
$q_n=\mathcal F(\theta[C]_{(n)})$.
His Remark~2 following Corollary~0.3, p.~3, says that
$(2g-2)P_0=K_C$ implies $q_n=0$ for all $n\ge1$.
The hypothesis holds for our hyperflex. On p.~4 he gives
$p_n\in\operatorname{CH}^n_{(n-1)}$ and $p_1=-\theta$.
Consequently the only codimension-two monomials in the remaining
generators are $p_1^2$ and $p_2=\eta$ (up to sign if the inverse
Fourier kernel is used). The class $\theta^2$ has nonzero
cohomology, while $\eta$ is homologically trivial and nonzero.
This proves (JAC.P12.1).
If $T\eta$ differed from a tautological class by an algebraically
trivial class, its $\theta^2$ coefficient would vanish by cohomology.
This would contradict the independence of $\eta,T\eta$ modulo
algebraic equivalence proved in JAC-P11.
\end{proof}

\paragraph{The remaining generation gate.}
The two-dimensional Betti realization of $N$ does not bound the
dimension of $\operatorname{CH}^2(N)$: every class in that Chow
group already has zero ordinary cycle class.
Finite-dimensionality of abelian Chow motives supplies nilpotence
for homologically trivial endomorphisms, not a bound on this Chow
group. The CM operation above produces a nonzero positive
one-dimensional $\mathbf Q(i)$-module, but does not prove that
the entire complementary Chow group has that dimension.
A generation theorem must also control a potential Abel--Jacobi
kernel. For comparison, Lilienfeldt--Shnidman,
\emph{Experiments with Ceresa classes of cyclic Fermat quotients},
author manuscript dated 29~November~2021, \S4.1, pp.~5--6,
labels the equality of Griffiths rank and the vanishing order of
the corresponding rank-two motivic $L$-function as Bloch's
conjectural \emph{Son of Recurring Fantasy}.
It is not a Chow-generation theorem.

\paragraph{Primary sources directly read.}
\url{https://www.jmilne.org/math/CourseNotes/AV.pdf},
\url{https://arxiv.org/pdf/math/0504448}, and
\url{https://math.huji.ac.il/~shnidman/Experiments_with_Ceresa_classes.pdf}.
The new deductions are the normalized CM isometry and positive
Griffiths plane in JAC-P11, and the explicit specialization of
the tautological generators in JAC-P12. The cited generation
theorem and Rosati positivity are used as established background.

